% ===================================================================
% =========================    摘要     =============================
% ===================================================================
%-
%-> 生成封面
%-

\maketitle% 生成中文封面
\MAKETITLE% 生成英文封面
%-
%-> 作者声明
%-
\makedeclaration% 生成声明页
%-
%-> 中文摘要
%-
\intobmk\chapter*{摘\quad 要}% 显示在书签但不显示在目录
\setcounter{page}{1}% 开始页码
\pagenumbering{Roman}% 页码符号

此处书写摘要。

文摘要包括中文摘要和英文摘要（Abstract）两部分。论文摘要应概括地反映出本论文的主要内容，说明本论文的主要研究目的、内容、方法、结论。要突出本论文的创造性成果或新见解，不宜使用公式、图表、表格或其他插图材料，不标注引用文献。中文摘要的字数由各学科群分会根据本分会涉及学科专业的特点提出具体要求。英文摘要与中文摘要内容应保持一致。留学生用其他语种撰写学位论文时，应有详细的中文摘要，字数由各学科群分会具体制定，建议一般不少于5000字。
摘要最后注明本文的关键词（3～5个）。关键词是为了文献标引和检索工作，从论文中选取出来，用以表示全文主题内容信息的单词或术语。关键词以显著的字符另起一行并隔行排列于摘要下方，左顶格，中文关键词间用中文逗号隔开。英文关键词应与中文关键词对应，首字母应大写，用英文逗号隔开。-- 《中国科学院大学研究生学位论文撰写规范指导意见（校发学位字〔2022〕40号）》



本文主要研究工作和贡献包括以下方面：

\begin{itemize}
    \item[(1)] xxx
    \item[(2)] xxx
    \item[(3)] xxx
    \item[(4)] xxx
\end{itemize}

\keywords{关键词 1，关键词 2，关键词 3}% 中文关键词
%-
%-> 英文摘要
%-
\intobmk\chapter*{Abstract}% 显示在书签但不显示在目录
\pagestyle{enfrontmatterstyle}%

English abstract goes here.

\begin{itemize}
    \item[(1)] xxx
    \item[(2)] xxx
    \item[(3)] xxx
    \item[(4)] xxx
\end{itemize}


\KEYWORDS{English Key Word 1, English Key Word 2, English Key Word 3}% 英文关键词实词首字母大写


\cleardoublepage\pagestyle{frontmatterstyle}%

%---------------------------------------------------------------------------%
